\chapter{O primeiro capitulo}
\label{cap:primeiro}

\epigraph{Uma epigrafe ambienta o capitulo em uma frase.}{--- \textup{Autor da frase}}

Este e o corpo do capitulo. Escreva em paragrafos, deixando que o LaTeX cuide da
quebra de linhas e da paginacao.

\section{Uma secao}
\label{sec:primeira}

Secoes organizam o capitulo. Termos que devem entrar no indice remissivo sao
marcados com \verb|\index|, assim: tipografia\index{tipografia} e
diagramacao\index{diagramacao}. O indice se monta sozinho ao fim do livro.

Para citar uma obra, use \verb|\cite| \cite{knuth1984} ou, dentro da frase,
\verb|\citeonline|, como em \citeonline{bringhurst2005}.

\subsection{Uma subsecao}

Ate tres niveis costumam bastar em um livro. Alem disso, o leitor se perde.

\section{Figuras e tabelas}
\label{sec:figuras-tabelas}

% Exemplo de figura. Descomente depois de colocar a imagem em figures/.
%
% \begin{figure}[htb]
%   \centering
%   \includegraphics[width=0.8\textwidth]{exemplo.png}
%   \caption{Descricao breve da figura}
%   \label{fig:exemplo}
%   \fonte{Elaborado pelo autor.}
% \end{figure}

\begin{table}[htb]
  \centering
  \caption{Uma tabela simples}
  \label{tab:exemplo}
  \begin{tabular}{lrr}
    \toprule
    \textbf{Item} & \textbf{Quantidade} & \textbf{Proporcao} \\
    \midrule
    Primeiro  & 120 & 40\% \\
    Segundo   &  90 & 30\% \\
    Terceiro  &  90 & 30\% \\
    \bottomrule
  \end{tabular}
  \fonte{Elaborado pelo autor.}
\end{table}

\begin{quote}
  Citacoes longas ganham um bloco proprio, recuado e sem aspas.
\end{quote}
